% C-35 -- methodology.tex anatomical/terminology consistency (redundancy_closing_audit_consolidation_2026-09-02.md, P-05 + P-24)
% Audit copy only -- live file is sections/methodology.tex.
%
% P-05: "femur" (proximal segment) was used where the rest of the paper (preamble.tex abstract,
% introduction.tex, and methodology.tex's own Physical Platform Construction subsection 47 lines
% below) consistently says "tibia" (distal segment) for the wheel-motor mounting point -- the
% wheel is at the distal end, so "tibia" is anatomically correct; "femur" was the outlier.
% P-24: "mecanum wheel" appeared exactly once in the same sentence, vs. "omnidirectional wheel"
% 14+ times elsewhere -- normalized for terminology consistency.

% --- methodology.tex, ~line 25 ---
This structure enables three degrees of freedom for each limb, maintaining the same design across the robot's four legs. Each limb includes three servomotors responsible for executing articulated movements, while an additional geared motor serves as the actuator for rolling motion. This motor is located at the distal end of the \revblue{tibia} and is directly coupled to an \revblue{omnidirectional} wheel, allowing smooth transitions between articulated and rolling modes.
